% fig_kappa_cascade.tex
% How kappa_n propagates through the 5-layer architecture
% Timeline plot showing kappa_n trajectory and the decisions it gates at each layer
\begin{tikzpicture}
\begin{groupplot}[
  group style={group size=1 by 2, vertical sep=0.45cm},
  thesis style,
  width=\columnwidth,
  xmin=0, xmax=200,
  xtick={0,40,80,120,160,200},
  grid=both,
  grid style={gray!20,line width=0.3pt},
  clip=true,
]

%% ── Panel 1: kappa_n trajectory ─────────────────────────────────────────────
\nextgroupplot[
  height=3.6cm,
  ylabel={$\kappan$ (condition number)},
  ymin=0, ymax=55,
  ytick={0,10,20,30,40,50},
  xticklabels={},
  legend pos=north east,
  legend style={font=\scriptsize,cells={anchor=west}},
]

%% kappa_n from LM estimator
\addplot[thick,thesisblue,mark=none,domain=0:200,samples=80]
  { (x<5) ? 50 :
    (x<60) ? max(6, 50*exp(-x/17)) :
    max(6, 7 + 0.045*(x-60)) };
\addlegendentry{$\kappan$ (L2 LM estimator)};

%% Threshold
\addplot[thesisred,thick,dotted,domain=0:200,samples=2] {30};
\node[thesisred,font=\tiny,anchor=west] at (axis cs:202,31.5) {30};

%% Layer-gate events
%% L3 adaptive law activates when kappa < 30
\draw[thesisgreen,thin,dashed] (axis cs:35,0) -- (axis cs:35,55)
  node[above,font=\tiny,thesisgreen,align=center]{L3\\active};

%% L4 CUEP gate
\draw[thesisblue,thin,dashed] (axis cs:35,0) -- (axis cs:35,55);
\node[thesisblue,font=\tiny,anchor=south] at (axis cs:50,30)
  {$\mathcal{K}(\kappan)$ narrows};

%% L5 veto lifted
\draw[thesisorange,thin,dashed] (axis cs:35,0) -- (axis cs:35,55);

%% Fill: kappa>30 (veto active, gray)
\addplot[thesisgray!25,fill=thesisgray!20,draw=none,domain=0:35,samples=2]
  {55} \closedcycle;
\node[thesisgray,font=\tiny,align=center] at (axis cs:17.5,45)
  {Veto\\active};

%% Fill: kappa<30 (decisions enabled, green tint)
\addplot[thesisgreen!10,fill=thesisgreen!8,draw=none,domain=35:200,samples=2]
  {30} \closedcycle;
\node[thesisgreen,font=\tiny,align=center] at (axis cs:115,22)
  {Gate open\\(decisions enabled)};

%% ── Panel 2: Layer output status over time ───────────────────────────────────
\nextgroupplot[
  height=3.0cm,
  ylabel={Layer output status},
  xlabel={Time after fault inception (ms)},
  ymin=0, ymax=5.5,
  ytick={1,2,3,4,5},
  yticklabels={L1,L2,L3,L4,L5},
]

%% Each layer: gray = inactive/waiting, coloured = active/valid
%% L1: always active (raw measurements)
\addplot[thesisblue!60,fill=thesisblue!25,draw=none]
  coordinates{(0,0.6)(200,0.6)(200,1.4)(0,1.4)} \closedcycle;
\node[font=\tiny,thesisblue] at (axis cs:100,1.0) {Active throughout};

%% L2: active after first LM iteration (~5ms)
\addplot[thesisorange!60,fill=thesisorange!20,draw=none]
  coordinates{(5,1.6)(200,1.6)(200,2.4)(5,2.4)} \closedcycle;
\addplot[thesisgray!30,fill=thesisgray!15,draw=none]
  coordinates{(0,1.6)(5,1.6)(5,2.4)(0,2.4)} \closedcycle;
\node[font=\tiny,thesisorange] at (axis cs:100,2.0) {Estimating $\hat\theta,\kappan$};

%% L3: active when kappa < 30 (from ~35ms)
\addplot[thesisgreen!60,fill=thesisgreen!20,draw=none]
  coordinates{(35,2.6)(200,2.6)(200,3.4)(35,3.4)} \closedcycle;
\addplot[thesisgray!30,fill=thesisgray!15,draw=none]
  coordinates{(0,2.6)(35,2.6)(35,3.4)(0,3.4)} \closedcycle;
\node[font=\tiny,thesisgreen] at (axis cs:115,3.0) {Adapting $k(t),\lambda(t)$};
\node[font=\tiny,thesisgray] at (axis cs:17.5,3.0) {Hold};

%% L4: CUEP check from ~35ms
\addplot[thesisred!50,fill=thesisred!15,draw=none]
  coordinates{(35,3.6)(200,3.6)(200,4.4)(35,4.4)} \closedcycle;
\addplot[thesisgray!30,fill=thesisgray!15,draw=none]
  coordinates{(0,3.6)(35,3.6)(35,4.4)(0,3.4)} \closedcycle;
\node[font=\tiny,thesisred] at (axis cs:115,4.0) {CUEP check: $V_0$ vs $V_{\rm cuep}$};

%% L5: decision from ~35ms
\addplot[thesisgray!50,fill=thesisgray!30,draw=none]
  coordinates{(35,4.6)(200,4.6)(200,5.4)(35,5.4)} \closedcycle;
\addplot[thesisgray!30,fill=thesisgray!15,draw=none]
  coordinates{(0,4.6)(35,4.6)(35,5.4)(0,5.4)} \closedcycle;
\node[font=\tiny,thesisgray] at (axis cs:115,5.0) {Trip / restrain};

%% veto boundary
\draw[thesisred,thin,dashed] (axis cs:35,0) -- (axis cs:35,5.5);

\end{groupplot}
\end{tikzpicture}
